\FloatBarrier
\section{本章小结}
\label{sec:tree-dp-chapter-summary}

本章的问题由局部函数组成：每个变量取有限个状态，整体成本是若干局部项之和，每一项只读取作用域中的少数变量，取值$+\infty$表示禁止的组合。同一份输入可以提出不同的输出要求：最小成本与一个最优配置、保留一组变量后的条件最小成本表、可行配置的个数，以及按局部权重之积定义的加权总和。枚举总能回答这些问题，但开销随变量数指数增长。

记录结构的几种图各有分工。原始图记录哪些变量共同出现，超图保留作用域的分组，因子图保留每个局部函数的身份；三者都不保存数值，原始图也不能恢复高阶规则，奇偶约束的例子说明三元规则不是成对规则的叠加。计算时始终使用完整的局部函数表。

链和树上之所以能精确递推，是因为删去依赖树的一条边就把变量分成两侧，两侧只通过一个接口变量相连。固定接口取值后，一侧的内部变量可以单独优化，结果写成以接口取值为索引的条件最小成本表；必须保留整张表，而不是一侧自己的最优值，因为接口最终取什么要等另一侧的成本也已知后才能决定。链上的递推、树上的消息和回溯表都是这一动作的具体形式，正确性来自有限集合上最小值的分步、移项和分离规则，开销对变量数是线性的。

依赖图含环时，单个变量不再能分隔两侧。消去一个变量会把它的邻居通过新函数耦合起来；四环算例中，消去$x_1$得到的$\psi_1(x_2,x_4)$不能拆成两张单变量表之和，接口必须同时保留两个变量。因此，一般图上的动态规划要先构建一族局部子问题，使每个局部项有处可放、接口足以概括内部变量、子问题的依赖可以排序；变量消元与树分解是描述这一构建的两种方式，分别在下一章建立并证明等价，数值计算统一放在第~\ref{chap:general-graph-dp-computation}章。
